\section{배경 및 관련 연구}\label{sec:related}

\subsection{SNN 효율 논거의 실체}\label{sec:premise}

SNN 문헌의 효율 논거는 통상 연산 단가표에서 출발한다. 45nm 32비트
부동소수점 기준으로 곱셈-누산은 4.6pJ, 덧셈은 0.9pJ라는
단가표\cite{horowitz2014computing}가 그것으로, 활성값이 이진인 SNN은
곱셈을 덧셈으로 대체하므로 발화율 $R$과 타임스텝 $T$가 작으면
$E_{AC}\!\cdot\!T\!\cdot\!R < E_{MAC}$로 이득이라는 형태의 손익
조건이 도출된다\cite{yao2023spikedriven}. 그러나 이 전제를 상용
프로세서로 옮기면 두 가지 정정이 필요하다. 첫째, 단가표는 데이터패스
\emph{회로}의 에너지에 관한 것이다. 파이프라인화된 MCU에서 곱셈은
이미 단일 사이클 명령(예: Cortex-M7의 SMLAD)이므로, ``곱셈을
덧셈으로 바꾸는 것'' 자체는 시간으로는 거의 아무것도 벌지 못한다.
둘째, 문헌의 정밀 분석들이 수렴하는 결론은 이득의 원천이 연산
치환이 아니라 \emph{희소성}이라는 것이다. Dampfhoffer 등은 디지털
하드웨어 모델에서 VGG16급 SNN의 시냅스 연산이 전체 에너지의 약
1\%에 불과하며 지배 항은 메모리 접근임을 보이고, SNN이 우세해지는
조건을 스파이크
희소성으로 귀속시킨다\cite{dampfhoffer2023snns}. Davidson과
Furber는 뉴런당 약 1.7 스파이크를 손익분기로 제시하며 rate 코딩
SNN의 이득에 회의적이고\cite{davidson2021comparison}, 반대편에는
6--8배 이득을 추정하는 분석도 있다\cite{lemaire2022analytical}.
추정이 가정에 따라 무이득에서 8배까지 갈린다는 사실 자체가, 고정된
상용 타깃 위의 \emph{실측}이 필요하다는 본 논문의 방법론적 입장을
뒷받침한다. 본 논문의 실측은 이 결론의 시간 영역 대응물을
제공한다: 인오더 MCU에서는 메모리 정지 사이클이 곧 시간 비용이므로
같은 구조가 지연에서도 나타나리라 기대할 수 있고, 실제로 병목은
0-피연산자의 메모리 이동으로 실측되었으며(III장), 그 트래픽을
제거한 커널이 단독 5.7배 --- 종단 구성 기준 6.12배 --- 를 냈다(V장).

\subsection{전용 실리콘의 회수}\label{sec:silicon}

이벤트 구동 이득을 하드웨어로 회수하는 길은 뉴로모픽 프로세서가
대표한다. TrueNorth\cite{merolla2014truenorth},
SpiNNaker\cite{furber2014spinnaker},
Loihi\cite{davies2018loihi}와 그 2세대\cite{orchard2021loihi2}는
스파이크가 있을 때만 연산이 일어나는 실행 모델을
실리콘$\cdot$시스템 수준에서 구현했고,
최근의 Speck은 완전 비동기 SoC에서 소비 전력이 입력 이벤트양에
비례함을 보였다\cite{yao2024speck}. 본 연구는 이 노선의 대체가
아니라 보완이다: 전용 실리콘 없이, 이미 보급된 상용 MCU에서, 같은
비례성의 \emph{산술적} 몫이 얼마나 회수되는지를 묻는다.

\subsection{희소성 활용 추론 시스템}\label{sec:sparse-sys}

DNN 일반에서 0-피연산자를 건너뛰려는 시도는 가속기 설계에서
시작되었다. Cnvlutin은 0 활성값의 곱셈을 동적으로
생략했고\cite{albericio2016cnvlutin}, EIE는 압축 표현 위에서
FC 층의 제로 스킵을\cite{han2016eie}, SCNN은 비영 쌍만 곱셈기에
전달하는 데이터플로를\cite{parashar2017scnn} 구현했다 --- 모두 전용
실리콘을 전제한다. 범용 CPU에서는 정적 가중치 희소성을 CSR류
형식으로 소비하는 직접 희소 합성곱(3--7배)\cite{park2017faster}과,
ReLU 활성 희소성을 학습으로 유도해 소비하는
커널\cite{kurtz2020inducing}이
서버급(x86, 광폭 벡터, GB급 메모리)에서 이득을 보였다.
공간 희소 데이터 쪽에서는 활성 좌표 집합에 합성곱을 제한하는
submanifold 방식\cite{graham2018submanifold}과 이벤트 카메라 입력의
비동기 희소 합성곱\cite{messikommer2020event}이 있으나, 전자는 GPU
라이브러리이고 후자의 절감은 주로 연산량(FLOP) 수준에서 보고된다.
우리가 확인한 범위에서 이 계열 중 MCU급(Cortex-M, RAM $<$1MB)
장치를 다룬 것은 없다. 또한 SNN의 이진 활성은 이 계열과 결이 다른
이점을 준다: 값이 $\{0,S\}$뿐이므로 희소 소비가 값 적재 없는
\emph{위치 전용} gather가 되고(IV장 원리 1), im2col류의 값
실체화\cite{chellapilla2006high} 자체가 원리적으로 불필요해진다.

\subsection{MCU 위의 SNN과 본 연구의 위치}\label{sec:positioning}

SNN 소프트웨어 스택은 학습 쪽으로
성숙했다\cite{eshraghian2023training,fang2023spikingjelly,pehle2021norse}.
그러나 배포 경로는 뉴로모픽 백엔드를 지향하거나 이식 표준의
단계\cite{pedersen2024nir}에 있고, MCU 급 상용 하드웨어를 향한
실측 연구는 이제 축적되는 중이다. Barchi 등은 STM32 계열에서 순환
SNN의 근접 센서 추론을 보였고\cite{barchi2021shm}, 후속 연구에서는
발화 활동 자체를 모델 수준에서 조율해 지연과 에너지를
낮췄다\cite{barchi2024tuning} --- 스파이크
수를 \emph{줄이는} 접근으로, 스파이크당 비용을 줄이는 본 연구와
직교한다. Jurzec 등은 snnTorch 모델을 자체 C 런타임으로 내려 Cortex-M7
보드에 올렸으나 희소성은 지속 비활성 뉴런의 오프라인 프루닝으로
쓰였고 속도 수치는 데스크톱 기준이다\cite{jurzec2025compression}.
개념적으로 가장 가까운 것은 Liao 등의 뇌-기계 인터페이스
디코더로, 병렬 RISC-V(GAP9)에서 발화한 뉴런의 가중치 열만 DMA로
가져오는 ``sparse copy''를 구현했다\cite{liao2024snndecoder} ---
발화로 가중치 트래픽을 게이팅한다는 착상을 공유하지만, 대상이
FC/순환 디코더라 수용 영역 인덱싱$\cdot$im2col 문제가 없고, 전용
런타임이며, 정확성은 근사 수준(``무시할 손실'')으로 보고된다.
SpikeStream은 희소 SNN 추론을 위해 RISC-V ISA에 간접 스트림
확장을 추가해 4.4배를 얻었는데\cite{manoni2025spikestream}, 이는
역으로 상용 ISA 그대로는 어렵다는 방증으로 읽을 수 있다 --- 본
연구는 그 간접 스트리밍을 전치 레이아웃과 레지스터 블로킹의 순수
소프트웨어로 대신한다(IV장 원리 2).

표~\ref{tab:positioning}이 위치를 요약한다. 우리가 조사한 범위에서,
\textbf{상용 Cortex-M급 MCU에서, 표준 TinyML 런타임 안에서, 합성곱
SNN의 발화 희소성을 런타임에 소비하며 조밀 경로와의 비트 동일성을
보증한 실측 보고는 본 연구가 처음이다}.

\begin{table*}[t]
\centering
\caption{희소성 활용 추론 연구들 속에서의 위치. ($\circ$: 해당,
$\triangle$: 부분적, $\times$: 비해당, --: 무관)}
\label{tab:positioning}
\footnotesize
\setlength{\tabcolsep}{5pt}
\begin{tabular}{l l l l c c}
\hline
연구 & 실행 기반 & 희소성 활용 방식 & 합성곱 & 표준 런타임 & 비트 동일 보증 \\
\hline
Speck\cite{yao2024speck} & 비동기 뉴로모픽 SoC & 실리콘 이벤트 구동 & $\circ$ & -- & -- \\
SCNN\cite{parashar2017scnn} & 전용 가속기(ASIC) & 비영 쌍 전달 하드웨어 & $\circ$ & -- & -- \\
Kurtz 등\cite{kurtz2020inducing} & x86 서버(AVX) & 유도 ReLU 희소성, SW & $\circ$ & $\times$ & $\triangle$ \\
SpikeStream\cite{manoni2025spikestream} & RISC-V 클러스터 & 간접 스트림 ISA 확장 & $\circ$ & $\times$ & $\times$ \\
Liao 등\cite{liao2024snndecoder} & GAP9 병렬 RISC-V & 발화 열 DMA 게이팅 & $\times$(FC/순환) & $\times$ & $\times$ \\
Barchi 등\cite{barchi2024tuning} & STM32(Cortex-M) & 모델 수준 발화 감소 & $\times$(순환) & $\times$ & $\times$ \\
Jurzec 등\cite{jurzec2025compression} & STM32H747(M7) & 오프라인 프루닝 & $\triangle$ & $\times$ & $\times$ \\
\textbf{본 연구} & \textbf{STM32F746(M7)} & \textbf{런타임 이벤트 게이팅, SW}
& $\circ$(수용 영역 인덱스) & $\circ$(TFLM$+$CMSIS-NN) & $\circ$(체크섬) \\
\hline
\end{tabular}
\begin{flushleft}
\footnotesize 주: ``비트 동일 보증''은 조밀 정수 경로와의 비트 단위
등가를 기계 검증하는지의 여부다. Kurtz 등의 커널은 수학적으로
정확하나 대상이 부동소수점 서버 추론이라 본 표의 관점(정수 경로
등가 검증)과는 구별된다.
\end{flushleft}
\end{table*}

\subsection{본 파이프라인의 발화 밀도}\label{sec:sparsity}

희소 커널의 이득은 전적으로 발화 밀도에 걸려 있으므로, 대상
파이프라인의 밀도를 먼저 실측해 둔다. 선행
연구\cite{cha2026fusion}는 희소성 활용을 향후 과제로 남기며 발화율을
보고하지 않았기에, 저자는 같은 모델$\cdot$같은 정수 파이프라인의
호스트 구현으로 CIFAR-10 테스트셋 10{,}000장$\times$30스텝을 전수
측정했다: 비영 활성 비율은 입력 스파이크 47.7\%, 첫 합성곱
출력(LIF1) 4.1\%, 둘째 출력(LIF2) 3.4\%다. 은닉층은 문헌의 통상
관측대로 수 퍼센트대이나 입력층은 rate 코딩 탓에 예외적으로
조밀하다 --- rate 코딩이 시간 코딩 대비 약 7배 많은 스파이크를
쓴다는 비교 분석\cite{guo2021neural}과 정합하는 값이다. 이 비대칭은
IV장의 비용 모델이 예측할 레이어 간 이득 격차의 원인이며, 스파이크
수 자체를 줄이는 학습 측 기법\cite{sorbaro2020optimizing}이 본
커널과 직교로 결합될 수 있는 이유이기도 하다.
